\begin{abstract}
In uniform internally heated convection between isothermal plates at infinite Prandtl number, a model of
radiogenically heated mantles, convection lowers the fraction $F_B$ of the heat that leaves through the base.
How fast can $F_B$ fall as the flux Rayleigh number $R$ grows? The best unconditional lower bounds, obtained with
the background method, were $F_B\gtrsim R^{-2/3}$ for no-slip and $R^{-40/29}$ for stress-free plates. We prove that
every classical solution satisfies $\overline{F_B}\ge c\,(R\mathcal T\ln R\mathcal T)^{-1/3}$ (no-slip) and
$\overline{F_B}\ge c\,(R\mathcal T)^{-1/2}$ (stress-free), where $\mathcal T$ is the long-time maximum
temperature, and that $\mathcal T\le C(1+\ln|\Omega|)$ for every incompressible flow. At fixed aspect ratio the
exponents therefore improve to $1/3$, up to a logarithm, and $1/2$. The flux argument does not use the energy
identities of the background method. The minimum principle bounds the horizontally averaged convective flux by
the maximal vertical speed times the mean temperature. A Green's-function estimate gives
$|w|\le\frac12R\mathcal T\Phi(z)$, with $\Phi(z)=\frac{3\sqrt3}{4\pi}z^2\ln(1/z)+O(z^2)$ in two dimensions. An
integrating factor then turns the averaged heat balance into an explicit bound, whose exact large-$R$ asymptotics
we determine. The constants depend on two slab-correction norms that we prove finite and evaluate numerically,
with numerical sensitivity checks, not uniform error enclosures but without interval certification. The rate $R^{-1/3}$ had previously been obtained only under
an unproved profile hypothesis, which the finite-window grid tests do not establish. Two-dimensional
Fourier--Chebyshev simulations and Newton-computed steady rolls, validated against published benchmarks, satisfy
every inequality of the proof with wide margins but stay far above the bound: over the computed ranges, published
time-dependent simulations and our $F_B$-minimising steady rolls both give $F_B\approx1.3R^{-0.09}$. For the basally heated problem at
infinite Prandtl number, within the family of mirror-symmetric, layer-spanning rolls and the period range we
scanned, the wide-roll maximum of the Nusselt number over the roll period, which reaches $\mathrm{Ra}^{1/3}$ at
unit Prandtl number, is not detected above the sampled interval between $\mathrm{Ra}=10^5$ and $1.8\times10^5$. Above that, the largest maxima
found are narrow rolls following $\mathrm{Nu}\approx0.15\,\mathrm{Ra}^{0.298}$ up to $\mathrm{Ra}=10^8$, close to
the high-wavenumber $\mathrm{Ra}^{3/10}$ regime. This extends a trend previously seen at Prandtl numbers 10 and
100, and it motivates the search for three-dimensional, unsteady or multi-scale constructions of solutions with
$\mathrm{Nu}\sim\mathrm{Ra}^{1/3}$.
\end{abstract}
